\subsection{Pass 10968: plane-rotation $R$-parity; the flat-or-massive dichotomy}

% Builds on Pass 10967; cf. Kappl et al. arXiv:0812.2120 (mu = 0 from an unbroken R-symmetry).
Adding the $W$-invariant combinations of the $\mathbb{Z}_6^R\times\mathbb{Z}_3^R\times\mathbb{Z}_2^R$ plane rotations of the
$\mathbb{Z}_6$-II lattice, matter parity exists in $3$ of $128$ models; $\mathbb{Z}_6$II$_{23}$ alone has parity-preserving
FI-cancelling $D$-flat directions ($38$ extreme rays; witness invariant on all $1524$ orbifolder-allowed couplings;
selection rules reproduce the orbifolder's $144$ order-$5$ $u^cd^cd^c$ couplings exactly).  In the $38$ the exotic
masses $u\,u^c$, $d\,d^c$, $e\,e^c$ are forbidden to all orders.  Two certified larger vacua break the remnant: $68$
directions (exotics full rank by order $5$, Higgs rank $6$) and $23$ directions (exotics full rank, Higgs rank $5/6$
at all orders, no outside linear term), each with one parity element and an exact positive $D$-flat vector; neither
has $W|_S=0$.  A branch search reached $202$ sub-vacua with $W|_S\equiv0$ ($71$ fully flat): the exotic $d$ mass matrix
has rank $0$ in all $202$.  $\mathbb{Z}_{12}$-I: $14$ of $289$ SMs reach a parity-viable vacuum and fail the exotic
mass test ($5$ robustly, $9$ with tentative $R$ rules).  $\mathbb{Z}_2\times\mathbb{Z}_6$-I: $29$ SMs, $0$ survive.
